% ============================================
% LUC-15 — A Mulher que Não Voltou para Casa (edição revista R1075)
% ============================================
\fragtitle{luc}{L}{LUC}{15}{A Mulher que Não Voltou para Casa}
\markboth{LUC-15 --- A Mulher que Não Voltou para Casa}{}

\noindent \textit{Belo Horizonte --- 9 de janeiro de 2026, 21h40}

\noindent Chovia na rua da Assembleia. Lúcia chegou ao bar com o casaco úmido e o cabelo colado na testa. Ronaldo já estava lá, numa mesa ao fundo, com um copo de uísque pela metade e o celular virado para baixo. Ele não a viu entrar. Estava olhando para a janela, para a chuva, com a expressão de quem espera uma ligação que não vai chegar.

\noindent Ela podia ter ido embora. Era o que fazia, normalmente. Quando o assunto era Ronaldo, ela sempre tinha uma reunião, um laudo, uma urgência. Desta vez, sentou.

\noindent ``Você demorou'', ele disse.

\noindent ``Você também.''

\noindent O bar cheirava a cerveja, madeira molhada e fritura. Mas, quando ele se inclinou para beijar seu rosto, Lúcia sentiu outra coisa: terra. Terra úmida, com um fundo doce, como se alguém tivesse aberto uma cova ali, entre as mesas. Olhou para o chão. Piso de cimento, seco. Olhou para ele. Ronaldo usava o mesmo perfume de sempre, madeira e âmbar. O cheiro de terra não estava nele.

\noindent ``Você está bem?'' ele perguntou.

\noindent ``Estou.''

\noindent Ela não estava. Havia três noites que acordava às 3h14. Havia duas semanas que o nome do paciente 1.260 aparecia em tudo o que ela lia. Havia oito anos que não ficava tão perto de Ronaldo.

\noindent Eles se conheceram no arquivo, em 2018. Ele era o diretor novo, ela a psiquiatra que vinha consultar processos antigos. Durou seis meses. Ele terminou num sábado, por telefone, dizendo que não sabia amar duas pessoas ao mesmo tempo. Lúcia não perguntou quem era a outra. Na segunda-feira, foi trabalhar com o rosto intacto e o estômago em fogo. Na sexta, aceitou um caso em outro estado. Quando voltou, ele estava casado.

\noindent Agora ele girava o copo devagar, com o dedo indicador na borda. A mão dele era a mesma. Lúcia olhou para o pulso, para o relógio antigo que ele usava, para o fio de cabelo branco que não existia oito anos atrás. Sentiu uma vontade absurda de encostar a testa no ombro dele. Não fez.

\noindent ``Preciso do acesso ao diário'', ela disse.

\noindent Ronaldo parou de girar o copo.

\noindent ``Não.''

\noindent ``Não é um pedido.''

\noindent ``Eu sei que não é.'' Ele levantou os olhos. ``É por isso que eu estou dizendo não.''

\noindent Lúcia abriu a pasta e colocou sobre a mesa a fotocópia de uma página do prontuário. O nome de Joaquim Antônio Correia. A data de 1979. A assinatura do Dr. Heitor Oliveira.

\noindent ``Você sabia que o diário estava lá'', ela disse. ``Sabia desde janeiro. Não me contou.''

\noindent ``Eu não sabia o que era.''

\noindent ``Sabia o suficiente para me mandar o e-mail.''

\noindent Ele ficou em silêncio. A chuva batia na vidraça. Lúcia viu, no reflexo do vidro, um cão parado na calçada, sob o toldo, olhando para dentro do bar. O olho direito do cão era amarelo. Não o amarelo dos faróis. Um amarelo denso, parado, como uma moeda velha. O cão não se mexeu.

\noindent ``Eu preciso ler'', ela disse. ``Se eu não ler, alguém vai ler no meu lugar. E eu não quero que seja alguém que não saiba o que está fazendo.''

\noindent Ronaldo riu sem humor.

\noindent ``Você sempre soube dizer isso. \textit{Eu preciso}. Como se o mundo fosse um paciente e você a única plantonista.''

\noindent ``Não é justo.''

\noindent ``Não é.'' Ele olhou para a mão dela sobre a mesa. ``Mas é verdade.''

\noindent Lúcia aproximou a mão. Encostou os dedos nos dele. Ronaldo não puxou. Ficaram assim por três segundos, quatro, cinco. O bar inteiro parecia ter baixado o volume. Ela sentiu o cheiro de terra de novo, mais forte, e desta vez veio com uma imagem: uma vala aberta, a terra revirada, uma mão de criança apontando para o céu.

\noindent Ela retirou a mão.

\noindent ``Me dá o cartão'', disse.

\noindent Ronaldo abriu a carteira. Tirou um cartão de acesso, desses que abrem a sala de documentos restritos. Colocou sobre a mesa, ao lado do copo.

\noindent ``Se você entrar nessa sala'', ele disse, ``não vai voltar para casa a mesma.''

\noindent Lúcia pegou o cartão. Guardou na pasta. Levantou.

\noindent ``Eu não vou voltar para casa hoje'', ela disse.

\noindent Ele não respondeu. Quando ela passou pela porta, o cão na calçada já tinha ido embora. No lugar onde ele estava, havia uma marca escura no chão, como se alguém tivesse derramado água. A chuva não lavava.

> **↪ Rotas:** **LUC-02** · **LUC-16** · **DOC-27**
